% BHU PYQ -- Manually transcribed & error-corrected
% Paper: BCF-123 : Fundamentals of Accounting
% Exam: B.Com. (Hons.) F.M.M. Semester II Examination, 2023-24

\documentclass[11pt,a4paper]{article}
\usepackage[a4paper,top=2.5cm,bottom=2.5cm,left=2.8cm,right=2.8cm,headheight=14pt]{geometry}
\usepackage{amsmath,amssymb}
\usepackage[T1]{fontenc}
\usepackage[utf8]{inputenc}
\usepackage{lmodern,microtype,fancyhdr,enumitem,hyperref,lastpage,parskip,booktabs}

\hypersetup{
    colorlinks=true,
    urlcolor=black,
    linkcolor=black,
    pdftitle={BCF-123: Fundamentals of Accounting -- BHU B.Com. (Hons.) FMM Sem II 2023-24}
}

\pagestyle{fancy}
\fancyhf{}
\renewcommand{\headrulewidth}{0.4pt}
\renewcommand{\footrulewidth}{0.4pt}
\fancyhead[L]{\small   \textit{Banaras Hindu University}}
\fancyhead[R]{\small   \textit{BCF-123: Fundamentals of Accounting}}
\fancyfoot[C]{\small Page  \thepage\ of \pageref{LastPage}}


\newlist{parts}{enumerate}{2}
\setlist[parts,1]{label=(\alph*),leftmargin=*,itemsep=5pt,topsep=3pt}
\setlist[parts,2]{label=(\roman*),leftmargin=*,itemsep=3pt,topsep=2pt}

\newcommand{\pts}[1]{\hfill   \textit{[#1]}}


\begin{document}


\begin{center}
  {\large\bfseries BANARAS HINDU UNIVERSITY}\\[2pt]
  {\normalsize B.Com. (Hons.) F.M.M. --- Semester II Examination, 2023-24}\\[6pt]
  \rule{0.8\linewidth}{0.5pt}\\[6pt]
  {\large\bfseries Paper No. BCF-123: Fundamentals of Accounting}\\[4pt]
  {\normalsize Time: 3 Hours\hfill Maximum Marks: 70}
\end{center}

\rule{\linewidth}{0.4pt}

\smallskip



\noindent   \textit{Instructions: Answer all questions. All questions carry equal marks (14 marks each).}

\bigskip




\noindent   \textbf{Question 1.}
From the following particulars, prepare a Bank Reconciliation Statement as on 31st March, 2023 by starting with (a) Cash Book Balances, and (b) Pass book Balance:\pts{14}

\begin{parts}
  \item On 31st March, the cash book of the firm showed a bank balance of Rs.\ 12,000 (Debit balance) while its pass book showed a bank balance of Rs.\ 10,360 (credit balance).
  \item Cheques of Rs.\ 10,000 out of which worth Rs.\ 8,000 only were presented for payment.
  \item Cheques of Rs.\ 2,800 were deposited in the bank on 27th March, 2023 but had not been credited by bank.
  \item One cheque for Rs.\ 1,000 entered in the cashbook on March, 2023 was sent to bank for collection on 01-04-2023.
  \item A cheque received from Geetika for Rs.\ 800 was deposited in bank on 26th March, 2023 but was dishonoured and the information was received on 02-04-2023.
  \item The pass book showed bank charges of Rs.\ 40 debited by bank.
  \item One of the debtors deposited a sum of Rs.\ 1,000 directly in the bank account of the firm on 28th March, 2023 but the intimation in this respect was received from the bank on 3-4-2023.
\end{parts}

\centerline{   \textbf{OR}}

Why the distinction between revenue and capital expenditure is of great importance in accounting? State the consideration which would guide you in deciding whether a particular item should be regarded as revenue or capital expenditure.\pts{14}

\medskip\rule{\linewidth}{0.2pt}




\noindent   \textbf{Question 2.}
Pass journal entries for the following transactions in the books of X Ltd.\ of Noida, Uttar Pradesh assuming CGST @ 6\% and SGST @ 6\% are levied:\pts{14}

\begin{parts}
  \item Purchased goods for Rs.\ 6,00,000 from Sayeba \& Co.\ of Patna, Bihar.
  \item Purchased goods for Rs.\ 1,00,000 from Gaurav Store of Varanasi, Uttar Pradesh.
  \item Sold goods costing Rs.\ 1,60,000 to Ishika of Ranchi, Jharkhand at a profit of 25\% on cost less 10\% Trade Discount.
  \item Sold goods costing Rs.\ 5,00,000 to Shubham of Allahabad, Uttar Pradesh at a profit of 60\% on cost less 15\% Trade Discount against cheque which was deposited into the bank.
  \item Paid for Advertisement Rs.\ 16,000.
  \item Purchased a computer for office use for Rs.\ 60,000 and payment was made by cheque.
  \item Proprietor withdrew Rs.\ 20,000 for his personal use.
  \item Payment made of the balance amount of GST.
\end{parts}

\centerline{   \textbf{OR}}

What is the purpose of a bank reconciliation statement? Give its proforma. What are the primary differences caused by errors in a bank reconciliation statement?\pts{14}

\medskip\rule{\linewidth}{0.2pt}




\noindent   \textbf{Question 3.}
From the following Trial balance and additional information, prepare Trading and Profit \& Loss Account for the year ended 31 March, 2024 and Balance sheet as on that date:\pts{14}


\begin{center}\small
  
\begin{tabular}{lr|lr}
     oprule
   \textbf{Particulars} &   \textbf{Debit (Rs.)} &   \textbf{Particulars} &   \textbf{Credit (Rs.)} \\
    \midrule
    Opening Stock & 33,000 & Sales & 3,00,000 \\
    Purchases & 2,00,000 & Sales tax collected & 32,000 \\
    Carriage inward & 3,300 & Creditors & 15,000 \\
    Carriage outward & 5,300 & Capital & 1,50,000 \\
    Salaries & 16,000 & Rent Received (1/3 portion of building) & 2,500 \\
    Administrative Expense & 13,400 & Provision for Bad \& Doubtful Debts & 500 \\
    Return inward & 6,700 & \quad (01/04/2023) & \\
    Sales Tax Paid & 24,000 & & \\
    Debtors & 25,000 & & \\
    Bad Debts & 2,000 & & \\
    Cash in hand & 2,300 & & \\
    Cash in bank & 9,000 & & \\
    Building & 1,00,000 & & \\
    Furniture & 60,000 & & \\
    \midrule
   \textbf{Total} &   \textbf{5,00,000} &   \textbf{Total} &   \textbf{5,00,000} \\
    ottomrule
  \end{tabular}
\end{center}




\noindent   \textbf{Additional Information:}

\begin{parts}
  \item Closing Stock on 31 March, 2024 amounted to Rs.\ 28,000.
  \item One-third portion of the building was used for business operations, one-third for residential purposes and one-third was let-out.
  \item New Bad Debts to Rs.\ 5,000, provision for Bad and Doubtful Debts is @ 5\% and create @ 2\% for Discount on Debtors and Reserve 2\% on Creditors.
  \item Charge depreciation on furniture @ 10\% p.a.
\end{parts}

\centerline{   \textbf{OR}}

What do you mean by adjustment entries? Why are they necessary for preparing final accounts? Give suitable example.\pts{14}

\medskip\rule{\linewidth}{0.2pt}




\noindent   \textbf{Question 4.}
A and B are partners in a firm with equal profit sharing ratio. On 1st April, 2023, they admit C as a new partner and decide to readjust the balance sheet values for the purpose. The Balance Sheet of A and B on 31st March, 2023 was as under:\pts{14}


\begin{center}\small
  
\begin{tabular}{lr|lr}
     oprule
   \textbf{Liabilities} &   \textbf{Rs.} &   \textbf{Assets} &   \textbf{Rs.} \\
    \midrule
    General Reserve & 1,000 & Cash & 7,000 \\
    Creditors & 10,000 & Debtors & 15,000 \\
    Bills Payable & 10,000 & Stock & 14,000 \\
    A's Capital & 17,000 & Furniture & 4,000 \\
    B's Capital & 12,000 & Machinery & 10,000 \\
    \midrule
   \textbf{Total} &   \textbf{50,000} &   \textbf{Total} &   \textbf{50,000} \\
    ottomrule
  \end{tabular}
\end{center}




\noindent C was admitted on the following terms:

\begin{parts}
  \item The value of stock to be increased by 10\% and that of Furniture to be decreased by 5\%, correct value of Machinery is Rs.\ 8,500.
  \item Investment worth Rs.\ 3,000 not shown in the Balance sheet are to be taken into account.
  \item An item of Rs.\ 500 included in creditors is not likely to be claimed.
  \item 5\% Provision for Doubtful debts is to be created on Debtors.
  \item C brings Rs.\ 10,000 as capital and Rs.\ 8,000 as goodwill. A and B withdraw half of the goodwill amount. C will get 1/4 share in the future profits.
\end{parts}
Prepare necessary accounts and New Balance Sheet.

\centerline{   \textbf{OR}}

Elucidate Joint Life Policy. What accounting records are maintained in the books of Partnership firm at the time of admission of New partner?\pts{14}

\medskip\rule{\linewidth}{0.2pt}




\noindent   \textbf{Question 5.}
A, B and C carry on business in partnership and divide profit and loss in the ratio 3 : 4 : 1. C retires from the firm. A and B agree to continue the business. They agree to show all the assets and liabilities of the firm at the value fixed by all three partners. On 30th September, 2023, the date of C's retirement, the position of the firm was as follows:\pts{14}


\begin{center}\small
  
\begin{tabular}{lr|lr}
     oprule
   \textbf{Liabilities} &   \textbf{Rs.} &   \textbf{Assets} &   \textbf{Rs.} \\
    \midrule
    Sundry Creditors & 1,89,000 & Plant & 40,000 \\
    Loan & 70,000 & Fixtures & 5,000 \\
    Capital Accounts: & & Furniture & 2,500 \\
    \quad A: 30,000 & & Stock & 1,00,000 \\
    \quad B: 40,000 & & Debtors & 1,80,000 \\
    \quad C: 10,000 & 80,000 & Bills Receivable & 40,000 \\
    Profit and Loss & 30,000 & Cash & 500 \\
    & & Cash in Bank & 1,000 \\
    \midrule
   \textbf{Total} &   \textbf{3,69,000} &   \textbf{Total} &   \textbf{3,69,000} \\
    ottomrule
  \end{tabular}
\end{center}




\noindent Assets were valued as under: Plant Rs.\ 30,500, Fixtures Rs.\ 3,500, Furniture Rs.\ 1,500, Stock at 20\% less, Sundry Debtors Rs.\ 1,50,000. A provision for Doubtful Bills Rs.\ 10,000 to be created.

\begin{parts}
  \item Prepare Capital Accounts of A, B and C after giving effect to the above adjustments.
  \item Prepare Balance Sheet showing the position of the new firm of A and B.
\end{parts}

\centerline{   \textbf{OR}}

Enumerate the various matters that need adjustments at the time of retirement of a partner.\pts{14}

\vfill
\rule{\linewidth}{0.4pt}

\begin{center}\small
\href{https://bhu-exam.vercel.app/}{Click here to visit our website}\\[4pt]
\href{https://me-aryan.vercel.app/}{Click here to visit our contributor}\\[4pt]
\href{https://quashan-me.vercel.app/}{Click here to visit our contributor}
\end{center}

\end{document}